\section{AHSME 1952}
        \begin{enumerate}
        	\item (\textit{AHSME 1952}) If the radius of a circle is a rational number, its area is given by a number which is:


 $\textbf{(A)\ } \text{rational}  \qquad \textbf{(B)\ } \text{irrational} \qquad \textbf{(C)\ } \text{integral} \qquad \textbf{(D)\ } \text{a perfect square }\qquad \textbf{(E)\ } \text{none of these}$


 

	\item (\textit{AHSME 1952}) Two high school classes took the same test. One class of $20$ students made an average grade of $80\%$; the other class of $30$ students made an average grade of $70\%$. The average grade for all students in both classes is:


 $\textbf{(A)}\ 75\%\qquad \textbf{(B)}\ 74\%\qquad \textbf{(C)}\ 72\%\qquad \textbf{(D)}\ 77\%\qquad \textbf{(E)\ }\text{none of these}$


 

	\item (\textit{AHSME 1952}) The expression $a^3-a^{-3}$ equals: 


 $\textbf{(A) \ }\left(a-\frac{1}{a}\right)\left(a^2+1+\frac{1}{a^2}\right) \qquad \textbf{(B) \ }\left(\frac{1}{a}-a\right)\left(a^2-1+\frac{1}{a^2}\right) \qquad \textbf{(C) \ }\left(a-\frac{1}{a}\right)\left(a^2-2+\frac{1}{a^2}\right) \qquad$
$\textbf{(D) \ }\left(\frac{1}{a}-a\right)\left(\frac{1}{a^2}+1+a^2\right) \qquad \textbf{(E) \ }\text{none of these}$


 

	\item (\textit{AHSME 1952}) The cost $C$ of sending a parcel post package weighing $P$ pounds, $P$ an integer, is $10$ cents for the first pound and $3$ cents for each additional pound. The formula for the cost is:


 $\textbf{(A) \ }C=10+3P \qquad \textbf{(B) \ }C=10P+3 \qquad \textbf{(C) \ }C=10+3(P-1) \qquad$


 $\textbf{(D) \ }C=9+3P \qquad \textbf{(E) \ }C=10P-7$


 

	\item (\textit{AHSME 1952}) The points $(6,12)$ and $(0,-6)$ are connected by a straight line. Another point on this line is:


 $\textbf{(A) \ }(3,3)  \qquad \textbf{(B) \ }(2,1) \qquad \textbf{(C) \ }(7,16) \qquad \textbf{(D) \ }(-1,-4) \qquad \textbf{(E) \ }(-3,-8)$


 

	\item (\textit{AHSME 1952}) The difference of the roots of $x^2-7x-9=0$ is:


 $\textbf{(A) \ }+7  \qquad \textbf{(B) \ }+\frac{7}{2} \qquad \textbf{(C) \ }+9 \qquad \textbf{(D) \ }2\sqrt{85} \qquad \textbf{(E) \ }\sqrt{85}$


 

	\item (\textit{AHSME 1952}) When simplified, $(x^{-1}+y^{-1})^{-1}$ is equal to:


 $\textbf{(A) \ }x+y  \qquad \textbf{(B) \ }\frac{xy}{x+y} \qquad \textbf{(C) \ }xy \qquad \textbf{(D) \ }\frac{1}{xy} \qquad \textbf{(E) \ }\frac{x+y}{xy}$


 

	\item (\textit{AHSME 1952}) Two equal circles in the same plane cannot have the following number of common tangents.


 $\textbf{(A) \ }1  \qquad \textbf{(B) \ }2 \qquad \textbf{(C) \ }3 \qquad \textbf{(D) \ }4 \qquad \textbf{(E) \ }\text{none of these}$


 

	\item (\textit{AHSME 1952}) If $m=\frac{cab}{a-b}$, then $b$ equals:


 $\textbf{(A) \ }\frac{m(a-b)}{ca}  \qquad \textbf{(B) \ }\frac{cab-ma}{-m} \qquad \textbf{(C) \ }\frac{1}{1+c} \qquad \textbf{(D) \ }\frac{ma}{m+ca} \qquad \textbf{(E) \ }\frac{m+ca}{ma}$


 

	\item (\textit{AHSME 1952}) An automobile went up a hill at a speed of $10$ miles an hour and down the same distance at a speed of $20$ miles an hour. The average speed for the round trip was:


 $\textbf{(A) \ }12\frac{1}{2}\text{mph}  \qquad \textbf{(B) \ }13\frac{1}{3}\text{mph} \qquad \textbf{(C) \ }14\frac{1}{2}\text{mph} \qquad \textbf{(D) \ }15\text{mph} \qquad \textbf{(E) \ }\text{none of these}$


 

	\item (\textit{AHSME 1952}) If $y=f(x)=\frac{x+2}{x-1}$, then it is incorrect to say:


 $\textbf{(A)\ }x=\frac{y+2}{y-1}\qquad\textbf{(B)\ }f(0)=-2\qquad\textbf{(C)\ }f(1)=0\qquad$


 $\textbf{(D)\ }f(-2)=0\qquad\textbf{(E)\ }f(y)=x$


 

	\item (\textit{AHSME 1952}) The sum to infinity of the terms of an infinite geometric progression is $6$. The sum of the first two terms is $4\frac{1}{2}$. The first term of the progression is:


 $\textbf{(A) \ }3 \text{ or } 1\frac{1}{2}  \qquad \textbf{(B) \ }1 \qquad \textbf{(C) \ }2\frac{1}{2} \qquad \textbf{(D) \ }6 \qquad \textbf{(E) \ }9\text{ or }3$


 

	\item (\textit{AHSME 1952}) The function $x^2+px+q$ with $p$ and $q$ greater than zero has its minimum value when:


 $\textbf{(A) \ }x=-p  \qquad \textbf{(B) \ }x=\frac{p}{2} \qquad \textbf{(C) \ }x=-2p \qquad \textbf{(D) \ }x=\frac{p^2}{4q} \qquad$


 $\textbf{(E) \ }x=\frac{-p}{2}$


 

	\item (\textit{AHSME 1952}) A house and store were sold for $\textdollar 12,000$ each. The house was sold at a loss of $20\%$ of the cost, and the store at a gain of $20\%$ of the cost. The entire transaction resulted in:


 $\textbf{(A) \ }\text{no loss or gain}  \qquad \textbf{(B) \ }\text{loss of }\textdollar 1000 \qquad \textbf{(C) \ }\text{gain of }\textdollar 1000 \qquad \textbf{(D) \ }\text{gain of }\textdollar 2000 \qquad \textbf{(E) \ }\text{none of these}$


 

	\item (\textit{AHSME 1952}) The sides of a triangle are in the ratio $6:8:9$. Then:


 $\textbf{(A) \ }\text{the triangle is obtuse}$


 $\textbf{(B) \ }\text{the angles are in the ratio }6:8:9$


 $\textbf{(C) \ }\text{the triangle is acute}$


 $\textbf{(D) \ }\text{the angle opposite the largest side is double the angle opposite the smallest side}$


 $\textbf{(E) \ }\text{none of these}$


 

	\item (\textit{AHSME 1952}) If the base of a rectangle is increased by $10\%$ and the area is unchanged, then the altitude is decreased by:


 $\textbf{(A) \ }9\%  \qquad \textbf{(B) \ }10\% \qquad \textbf{(C) \ }11\% \qquad \textbf{(D) \ }11\frac{1}{9}\% \qquad \textbf{(E) \ }9\frac{1}{11}\%$


 

	\item (\textit{AHSME 1952}) A merchant bought some goods at a discount of $20\%$ of the list price. He wants to mark them at such a price that he can give a discount of $20\%$ of the marked price and still make a profit of $20\%$ of the selling price. The percent of the list price at which he should mark them is:


 $\textbf{(A) \ }20  \qquad \textbf{(B) \ }100 \qquad \textbf{(C) \ }125 \qquad \textbf{(D) \ }80 \qquad \textbf{(E) \ }120$


 

	\item (\textit{AHSME 1952}) $\log p+\log q=\log(p+q)$ only if:


 $\textbf{(A) \ }p=q=\text{zero}  \qquad \textbf{(B) \ }p=\frac{q^2}{1-q} \qquad \textbf{(C) \ }p=q=1 \qquad$


 $\textbf{(D) \ }p=\frac{q}{q-1} \qquad \textbf{(E) \ }p=\frac{q}{q+1}$


 

	\item (\textit{AHSME 1952}) Angle $B$ of triangle $ABC$ is trisected by $BD$ and $BE$ which meet $AC$ at $D$ and $E$ respectively. Then:


 $\textbf{(A) \ }\frac{AD}{EC}=\frac{AE}{DC}  \qquad \textbf{(B) \ }\frac{AD}{EC}=\frac{AB}{BC} \qquad \textbf{(C) \ }\frac{AD}{EC}=\frac{BD}{BE} \qquad$


 $\textbf{(D) \ }\frac{AD}{EC}=\frac{(AB)(BD)}{(BE)(BC)} \qquad \textbf{(E) \ }\frac{AD}{EC}=\frac{(AE)(BD)}{(DC)(BE)}$


 

	\item (\textit{AHSME 1952}) If $\frac{x}{y}=\frac{3}{4}$, then the incorrect expression in the following is:


 $\textbf{(A) \ }\frac{x+y}{y}=\frac{7}{4}  \qquad \textbf{(B) \ }\frac{y}{y-x}=\frac{4}{1} \qquad \textbf{(C) \ }\frac{x+2y}{x}=\frac{11}{3} \qquad$


 $\textbf{(D) \ }\frac{x}{2y}=\frac{3}{8} \qquad \textbf{(E) \ }\frac{x-y}{y}=\frac{1}{4}$


 

	\item (\textit{AHSME 1952}) The sides of a regular polygon of $n$ sides, $n>4$, are extended to form a star. The number of degrees at each point of the star is:


 $\textbf{(A) \ }\frac{360}{n}  \qquad \textbf{(B) \ }\frac{(n-4)180}{n} \qquad \textbf{(C) \ }\frac{(n-2)180}{n} \qquad$


 $\textbf{(D) \ }180-\frac{90}{n} \qquad \textbf{(E) \ }\frac{180}{n}$


 

	\item (\textit{AHSME 1952}) On hypotenuse $AB$ of a right triangle $ABC$ a second right triangle $ABD$ is constructed with hypotenuse $AB$. If $\overline{BC}=1$, $\overline{AC}=b$, and $\overline{AD}=2$, then $\overline{BD}$ equals:


 $\textbf{(A) \ }\sqrt{b^2+1}  \qquad \textbf{(B) \ }\sqrt{b^2-3} \qquad \textbf{(C) \ }\sqrt{b^2+1}+2 \qquad$


 $\textbf{(D) \ }b^2+5 \qquad \textbf{(E) \ }\sqrt{b^2+3}$


 

	\item (\textit{AHSME 1952}) If $\frac{x^2-bx}{ax-c}=\frac{m-1}{m+1}$ has roots which are numerically equal but of opposite signs, the value of $m$ must be:


 $\textbf{(A) \ }\frac{a-b}{a+b} \qquad \textbf{(B) \ }\frac{a+b}{a-b} \qquad \textbf{(C) \ }c \qquad \textbf{(D) \ }\frac{1}{c} \qquad \textbf{(E) \ }1$


 

	\item (\textit{AHSME 1952}) In the figure, it is given that angle $C = 90^{\circ}$,$\overline{AD} = \overline{DB}$,$DE \perp AB$, $\overline{AB} = 20$, and $\overline{AC} = 12$. The area of quadrilateral $ADEC$ is: 

                    \begin{figure}[H]
                        \centering
                        \includegraphics[width=0.5\textwidth]{Problem/AHSME_I_1952_24_Problem_Figure_1.png}
                    \end{figure}
                    


 $\textbf{(A)}\ 75\qquad\textbf{(B)}\ 58\frac{1}{2}\qquad\textbf{(C)}\ 48\qquad\textbf{(D)}\ 37\frac{1}{2}\qquad\textbf{(E)}\ \text{none of these}$


 

	\item (\textit{AHSME 1952}) A powderman set a fuse for a blast to take place in $30$ seconds. He ran away at a rate of $8$ yards per second. Sound travels at the rate of $1080$ feet per second. When the powderman heard the blast, he had run approximately:


 $\textbf{(A)}\ \text{200 yd.}\qquad\textbf{(B)}\ \text{352 yd.}\qquad\textbf{(C)}\ \text{300 yd.}\qquad\textbf{(D)}\ \text{245 yd.}\qquad\textbf{(E)}\ \text{512 yd.}$


 

	\item (\textit{AHSME 1952}) If $\left(r+\frac1r\right)^2=3$, then $r^3+\frac1{r^3}$ equals


 $\textbf{(A)}\ 1\qquad\textbf{(B)}\ 2\qquad\textbf{(C)}\ 0\qquad\textbf{(D)}\ 3\qquad\textbf{(E)}\ 6$


 

	\item (\textit{AHSME 1952}) The ratio of the perimeter of an equilateral triangle having an altitude equal to the radius of a circle, to the perimeter of an equilateral triangle inscribed in the circle is:


 $\textbf{(A)}\ 1:2\qquad\textbf{(B)}\ 1:3\qquad\textbf{(C)}\ 1:\sqrt3\qquad\textbf{(D)}\ \sqrt3:2 \qquad\textbf{(E)}\ 2:3$


 

	\item (\textit{AHSME 1952}) In the table shown, the formula relating $x$ and $y$ is:


 $\begin{tabular}{|l|l|l|l|l|l|} \hline x & 1 & 2 & 3 & 4 & 5 \\ \hline y & 3 & 7 & 13 & 21 &  31 \\ \hline \end{tabular}$


 

$\textbf{(A)}\ y=4x-1  \qquad \textbf{(B)}\ y=x^{3}-x^{2}+x+2 \qquad \textbf{(C)}\ y=x^2+x+1 \\ \qquad \textbf{(D)}\ y=(x^2+x+1)(x-1) \qquad \textbf{(E)}\ \text{None of these}$


 

	\item (\textit{AHSME 1952}) In a circle of radius $5$ units, $CD$ and $AB$ are perpendicular diameters. A chord $CH$ cutting $AB$ at $K$ is $8$ units long.  The diameter $AB$ is divided into two segments whose dimensions are:


 $\textbf{(A)}\ 1.25, 8.75  \qquad \textbf{(B)}\ 2.75, 7.25\qquad \textbf{(C)}\ 2,8 \qquad \textbf{(D)} \ 4,6 \qquad \textbf{(E)} \text{None of these}$


 

	\item (\textit{AHSME 1952}) When the sum of the first ten terms of an arithmetic progression is four times the sum of the first five terms, the ratio of the first term to the common difference is: 


 $\textbf{(A)}\ 1: 2 \qquad \textbf{(B)}\ 2: 1 \qquad \textbf{(C)}\ 1: 4 \qquad \textbf{(D)}\ 4: 1 \qquad \textbf{(E)}\ 1: 1$


 

	\item (\textit{AHSME 1952}) Given $12$ points in a plane no three of which are collinear, the number of lines they determine is: 


 $\textbf{(A)}\ 24 \qquad \textbf{(B)}\ 54 \qquad \textbf{(C)}\ 120 \qquad \textbf{(D)}\ 66 \qquad \textbf{(E)}\ \text{none of these}$


 

	\item (\textit{AHSME 1952}) $K$ takes $30$ minutes less time than $M$ to travel a distance of $30$ miles. $K$ travels $\frac {1}{3}$ mile per hour faster than $M$. If $x$ is $K$'s rate of speed in miles per hours, then $K$'s time for the distance is: 


 $\textbf{(A)}\ \dfrac{x + \frac {1}{3}}{30} \qquad \textbf{(B)}\ \dfrac{x - \frac {1}{3}}{30} \qquad \textbf{(C)}\ \frac{30}{x+\frac{1}{3}}\qquad \textbf{(D)}\ \frac{30}{x}\qquad \textbf{(E)}\ \frac{x}{30}$


 

	\item (\textit{AHSME 1952}) A circle and a square have the same perimeter. Then: 


 $\text{(A) their areas are equal}\qquad\\ \text{(B) the area of the circle is the greater} \qquad\\ \text{(C) the area of the square is the greater} \qquad\\ \text{(D) the area of the circle is } \pi \text{ times the area of the square}\qquad\\ \text{(E) none of these}$


 

	\item (\textit{AHSME 1952}) The price of an article was increased $p\%$. Later the new price was decreased $p\%$. If the last price was one dollar, the original price was: 


 $\textbf{(A)}\ \frac{1-p^2}{200}\qquad \textbf{(B)}\ \frac{\sqrt{1-p^2}}{100}\qquad \textbf{(C)}\ \text{one dollar}\qquad\\ \textbf{(D)}\ 1-\frac{p^2}{10000-p^2}\qquad \textbf{(E)}\ \frac{10000}{10000-p^2}$


 

	\item (\textit{AHSME 1952}) With a rational denominator, the expression $\frac {\sqrt {2}}{\sqrt {2} + \sqrt {3} - \sqrt {5}}$ is equivalent to: 


 $\textbf{(A)}\ \frac {3 + \sqrt {6} + \sqrt {15}}{6} \qquad \textbf{(B)}\ \frac {\sqrt {6} - 2 + \sqrt {10}}{6} \qquad \textbf{(C)}\ \frac{2+\sqrt{6}+\sqrt{10}}{10} \qquad\\ \textbf{(D)}\ \frac {2 + \sqrt {6} - \sqrt {10}}{6} \qquad \textbf{(E)}\ \text{none of these}$


 

	\item (\textit{AHSME 1952}) To be continuous at $x = - 1$, the value of $\frac {x^3 + 1}{x^2 - 1}$ is taken to be: 


 $\textbf{(A)}\ - 2 \qquad \textbf{(B)}\ 0 \qquad \textbf{(C)}\ \frac {3}{2} \qquad \textbf{(D)}\ \infty \qquad \textbf{(E)}\ -\frac{3}{2}$


 

	\item (\textit{AHSME 1952}) Two equal parallel chords are drawn $8$ inches apart in a circle of radius $8$ inches. The area of that part of the circle that lies between the chords is:


 $\textbf{(A)}\ 21\frac{1}{3}\pi-32\sqrt{3}\qquad \textbf{(B)}\ 32\sqrt{3}+21\frac{1}{3}\pi\qquad \textbf{(C)}\ 32\sqrt{3}+42\frac{2}{3}\pi \qquad\\ \textbf{(D)}\ 16\sqrt {3} + 42\frac {2}{3}\pi \qquad \textbf{(E)}\ 42\frac {2}{3}\pi$


 

	\item (\textit{AHSME 1952}) The area of a trapezoidal field is $1400$ square yards. Its altitude is $50$ yards. Find the two bases, if the number of yards in each base is an integer divisible by $8$. The number of solutions to this problem is: 


 $\textbf{(A)}\ \text{none} \qquad \textbf{(B)}\ \text{one} \qquad \textbf{(C)}\ \text{two} \qquad \textbf{(D)}\ \text{three} \qquad \textbf{(E)}\ \text{more than three}$


 

	\item (\textit{AHSME 1952}) If the perimeter of a rectangle is $p$ and its diagonal is $d$, the difference between the length and width of the rectangle is: 


 $\textbf{(A)}\ \frac {\sqrt {8d^2 - p^2}}{2} \qquad \textbf{(B)}\ \frac {\sqrt {8d^2 + p^2}}{2} \qquad \textbf{(C)}\ \frac{\sqrt{6d^2-p^2}}{2}\qquad\\ \textbf{(D)}\ \frac {\sqrt {6d^2 + p^2}}{2} \qquad \textbf{(E)}\ \frac {8d^2 - p^2}{4}$


 

	\item (\textit{AHSME 1952}) In order to draw a graph of $ax^2+bx+c$, a table of values was constructed. These values of the function for a set of equally spaced increasing values of $x$ were $3844, 3969, 4096, 4227, 4356, 4489, 4624$, and $4761$. The one which is incorrect is: 


 $\text{(A) } 4096 \qquad \text{(B) } 4356 \qquad \text{(C) } 4489 \qquad \text{(D) } 4761 \qquad \text{(E) } \text{none of these}$


 

	\item (\textit{AHSME 1952}) Increasing the radius of a cylinder by $6$ units increased the volume by $y$ cubic units. Increasing the altitude of the cylinder by $6$ units also increases the volume by $y$ cubic units. If the original altitude is $2$, then the original radius is: 


 $\text{(A) } 2 \qquad \text{(B) } 4 \qquad \text{(C) } 6 \qquad \text{(D) } 6\pi \qquad \text{(E) } 8$


 

	\item (\textit{AHSME 1952}) Let $D$ represent a repeating decimal. If $P$ denotes the $r$ figures of $D$ which do not repeat themselves, and $Q$ denotes the $s$ figures of $D$ which do repeat themselves, then the incorrect expression is:


 $\textbf{(A) } D = .PQQQ\ldots \qquad\\ \textbf{(B) } 10^rD = P.QQQ\ldots \\ \textbf{(C) } 10^{r + s}D = PQ.QQQ\ldots \qquad\\ \textbf{(D) } 10^r(10^s - 1)D = Q(P - 1) \\ \textbf{(E) } 10^r\cdot10^{2s}D = PQQ.QQQ\ldots$


 

	\item (\textit{AHSME 1952}) The diameter of a circle is divided into $n$ equal parts. On each part a semicircle is constructed. As $n$ becomes very large, the sum of the lengths of the arcs of the semicircles approaches a length: 


 $\textbf{(A) } \qquad$ equal to the semi-circumference of the original circle


 $\textbf{(B) } \qquad$ equal to the diameter of the original circle


 $\textbf{(C) }$ greater than the diameter, but less than the semi-circumference of the original circle


 $\textbf{(D) } \qquad$ that is infinite


 $\textbf{(E) }$ greater than the semi-circumference


 

	\item (\textit{AHSME 1952}) If an integer of two digits is $k$ times the sum of its digits, the number formed by interchanging the the digits is the sum of the digits multiplied by


 $\textbf{(A) \ } 9-k  \qquad \textbf{(B) \ } 10-k \qquad \textbf{(C) \ } 11-k \qquad \textbf{(D) \ } k-1 \qquad \textbf{(E) \ } k+1$


 

	\item (\textit{AHSME 1952}) If $a$ and $b$ are two unequal positive numbers, then: 


 $\text{(A) } \frac{2ab}{a+b}>\sqrt{ab}>\frac{a+b}{2}\qquad \text{(B) } \sqrt{ab}>\frac{2ab}{a+b}>\frac{a+b}{2} \\ \text{(C) } \frac{2ab}{a+b}>\frac{a+b}{2}>\sqrt{ab}\qquad \text{(D) } \frac{a+b}{2}>\frac{2ab}{a+b}>\sqrt{ab} \\ \text{(E) } \frac {a + b}{2} > \sqrt {ab} > \frac {2ab}{a + b}$


 

	\item (\textit{AHSME 1952}) The base of a new rectangle equals the sum of the diagonal and the greater side of a given rectangle, while the altitude of the new rectangle equals the difference of the diagonal and the greater side of the given rectangle. The area of the new rectangle is: 


 $\text{(A) greater than the area of the given rectangle} \quad\\ \text{(B) equal to the area of the given rectangle}  \quad\\ \text{(C) equal to the area of a square with its side equal to the smaller side of the given rectangle}  \quad\\ \text{(D) equal to the area of a square with its side equal to the greater side of the given rectangle}  \quad\\ \text{(E) equal to the area of a rectangle whose dimensions are the diagonal and the shorter side of the given rectangle}$


 

	\item (\textit{AHSME 1952}) In the set of equations $z^x = y^{2x},\quad  2^z = 2\cdot4^x, \quad x + y + z = 16$, the integral roots in the order $x,y,z$ are: 


 $\textbf{(A) } 3,4,9 \qquad \textbf{(B) } 9,-5,-12 \qquad \textbf{(C) } 12,-5,9 \qquad \textbf{(D) } 4,3,9 \qquad \textbf{(E) } 4,9,3$


 

	\item (\textit{AHSME 1952}) Two cyclists, $k$ miles apart, and starting at the same time, would be together in $r$ hours if they traveled in the same direction, but would pass each other in $t$ hours if they traveled in opposite directions. The ratio of the speed of the faster cyclist to that of the slower is: 


 $\text{(A) } \frac {r + t}{r - t} \qquad \text{(B) } \frac {r}{r - t} \qquad \text{(C) } \frac {r + t}{r} \qquad \text{(D) } \frac{r}{t}\qquad \text{(E) } \frac{r+k}{t-k}$


 

	\item (\textit{AHSME 1952}) 
                    \begin{figure}[H]
                        \centering
                        \includegraphics[width=0.5\textwidth]{Problem/AHSME_I_1952_49_Problem_Figure_1.png}
                    \end{figure}
                    


 In the figure, $\overline{CD}$, $\overline{AE}$ and $\overline{BF}$ are one-third of their respective sides. It follows that $\overline{AN_2}: \overline{N_2N_1}: \overline{N_1D} = 3: 3: 1$, and similarly for lines BE and CF. Then the area of triangle $N_1N_2N_3$ is: 


 $\text{(A) } \frac {1}{10} \triangle ABC \qquad \text{(B) } \frac {1}{9} \triangle ABC \qquad \text{(C) } \frac{1}{7}\triangle ABC\qquad \text{(D) } \frac{1}{6}\triangle ABC\qquad \text{(E) } \text{none of these}$


 

	\item (\textit{AHSME 1952}) A line initially 1 inch long grows according to the following law, where the first term is the initial length. 
$1+\frac{1}{4}\sqrt{2}+\frac{1}{4}+\frac{1}{16}\sqrt{2}+\frac{1}{16}+\frac{1}{64}\sqrt{2}+\frac{1}{64}+\cdots$


 If the growth process continues forever, the limit of the length of the line is: 


 $\textbf{(A) } \infty\qquad \textbf{(B) } \frac{4}{3}\qquad \textbf{(C) } \frac{8}{3}\qquad \textbf{(D) } \frac{1}{3}(4+\sqrt{2})\qquad \textbf{(E) } \frac{2}{3}(4+\sqrt{2})$


 

\end{enumerate}